\documentclass[12pt,a4paper]{report}

% ============================================================
% PACKAGES
% ============================================================
\usepackage[utf8]{inputenc}
\usepackage[T1]{fontenc}
\usepackage[vietnamese,english]{babel}
\usepackage{lmodern}
\usepackage[top=2.5cm, bottom=2.5cm, left=3cm, right=2.5cm]{geometry}
\usepackage{amsmath, amssymb, amsfonts}
\usepackage{graphicx}
\usepackage{tikz}
\usetikzlibrary{shapes.geometric, arrows.meta, positioning, fit, calc, backgrounds}
\usepackage{algorithm}
\usepackage{algpseudocode}
\usepackage{booktabs}
\usepackage{hyperref}
\usepackage{cleveref}
\usepackage{enumitem}
\usepackage{caption}
\usepackage{xcolor}
\usepackage{setspace}
\usepackage{fancyhdr}

\onehalfspacing

\definecolor{corecolor}{HTML}{2E86AB}
\definecolor{ext1color}{HTML}{A23B72}
\definecolor{ext2color}{HTML}{F18F01}
\definecolor{optcolor}{HTML}{6B8E23}

\pagestyle{fancy}
\fancyhf{}
\fancyhead[L]{\small\leftmark}
\fancyhead[R]{\small\thepage}
\renewcommand{\headrulewidth}{0.4pt}

\hypersetup{colorlinks=true, linkcolor=corecolor, citecolor=corecolor, urlcolor=corecolor}

% ============================================================
\title{
    \vspace{-1cm}
    \textbf{Đề Cương Đồ Án Tốt Nghiệp} \\[0.5cm]
    \Large\textbf{Hệ Thống Nhận Dạng Từ Khóa Thoại Few-Shot Open-Set \\
    với Phân Lớp Prototype Kháng Nhiễu \\
    và Suy Luận Streaming Thời Gian Thực} \\[1cm]
    \large Khoa Công nghệ Thông tin \\
    \large Năm học 2025--2026
}
\author{Họ tên sinh viên \\ Giảng viên hướng dẫn: [Tên GVHD]}
\date{\today}

\begin{document}

\selectlanguage{vietnamese}
\maketitle
\thispagestyle{empty}

% ============================================================
\begin{abstract}
Đề cương này trình bày thiết kế và kế hoạch triển khai một hệ thống Nhận dạng
Từ khóa Thoại (Keyword Spotting -- KWS) theo phương pháp Few-Shot Open-Set, giải
quyết ba thách thức chính: (1) cho phép người dùng tùy chỉnh từ khóa chỉ với
3--5 mẫu âm thanh mà không cần huấn luyện lại mô hình, (2) hoạt động ổn định
trong môi trường có tiếng ồn, và (3) suy luận trực tuyến theo thời gian thực.

Hệ thống sử dụng mạng tích chập phân tách theo chiều sâu (DSCNN) được huấn luyện
bằng Triplet Loss trên bộ dữ liệu MSWC để học các biểu diễn âm thanh phân biệt.
Chúng tôi đề xuất phương pháp phân lớp dựa trên khoảng cách L2 trực tiếp (Direct
L2 Distance), loại bỏ bước chuẩn hóa xác suất nhân tạo, kết hợp với module khử
nhiễu âm thanh và bộ suy luận streaming tích hợp VAD (Voice Activity Detection).
Module xác thực người nói (Speaker Verification) sử dụng ECAPA-TDNN được cung cấp
như tính năng tùy chọn để cá nhân hóa hệ thống.

\textbf{Từ khóa:} Nhận dạng Từ khóa Thoại, Học Few-Shot, Nhận dạng Open-Set,
Mạng Prototypical, DSCNN, Suy luận Streaming, Phát hiện Hoạt động Giọng nói.
\end{abstract}

\tableofcontents
\newpage

% ============================================================
\chapter{Giới Thiệu}

\section{Bối cảnh và Động lực}

Nhận dạng Từ khóa Thoại (Keyword Spotting -- KWS) là bài toán cơ bản trong xử lý
tiếng nói, yêu cầu hệ thống tự động phát hiện các từ khóa mục tiêu được định nghĩa
trước trong luồng âm thanh liên tục. Công nghệ này là nền tảng cho nhiều ứng dụng
thực tế:

\begin{itemize}
    \item Trợ lý ảo kích hoạt bằng giọng nói (Amazon Alexa, Google Assistant)
    \item Hệ thống điều khiển rảnh tay trong ô tô
    \item Thiết bị nhà thông minh với từ đánh thức (``Hey Siri'', ``OK Google'')
\end{itemize}

Các hệ thống KWS truyền thống được huấn luyện trên bộ từ vựng cố định, lớn và
yêu cầu huấn luyện lại toàn bộ khi người dùng muốn thêm từ khóa mới. Hạn chế
này đặc biệt nghiêm trọng khi:

\begin{enumerate}
    \item \textbf{Nhu cầu cá nhân hóa:} Người dùng mong muốn thiết bị phản hồi
          các lệnh giọng nói tùy chỉnh theo nhu cầu riêng.
    \item \textbf{Ràng buộc tài nguyên:} Việc huấn luyện lại trên thiết bị biên
          (edge device) tốn kém về tính toán và không thực tế cho các thiết bị
          có bộ nhớ và năng lực xử lý hạn chế.
\end{enumerate}

Phương pháp \textbf{Few-shot Learning} giải quyết vấn đề này bằng cách cho phép
hệ thống học từ khóa mới chỉ từ 3--5 mẫu, mà không cần huấn luyện lại mô hình.
Kết hợp với \textbf{Open-set Recognition} --- khả năng từ chối các từ không xác
định (không thuộc bất kỳ từ khóa nào đã đăng ký) --- tạo nên hệ thống KWS có
khả năng triển khai thực tế.

\section{Mục tiêu Nghiên cứu}

Đồ án này giải quyết các câu hỏi nghiên cứu sau:

\begin{enumerate}
    \item Làm thế nào xây dựng hệ thống KWS có thể học từ khóa mới từ ít mẫu
          (few-shot) đồng thời từ chối đúng các từ lạ (open-set)?
    \item Phương pháp phân lớp dựa trên khoảng cách L2 trực tiếp có vượt trội
          so với phương pháp dựa trên xác suất không?
    \item Việc tích hợp khử nhiễu âm thanh ảnh hưởng thế nào đến độ chính xác
          KWS trong các điều kiện nhiễu khác nhau?
    \item Hệ thống có duy trì được độ chính xác chấp nhận được khi mở rộng từ
          offline sang suy luận streaming không?
\end{enumerate}

\section{Đóng góp Cốt lõi}

\begin{enumerate}
    \item \textbf{Đề xuất cơ chế phân lớp prototype dựa trên khoảng cách L2
          trực tiếp} cho bài toán KWS few-shot open-set, loại bỏ ràng buộc
          chuẩn hóa xác suất, cung cấp diễn giải ngưỡng trực quan như ``bán
          kính chấp nhận'' quanh mỗi prototype.
    \item \textbf{Đánh giá tác động của noise robustness} thông qua kết hợp
          noise augmentation khi huấn luyện và denoising preprocessing khi
          suy luận.
    \item \textbf{Mở rộng hệ thống từ offline sang streaming} với VAD và
          sliding-window detection.
    \item \textbf{(Tùy chọn) Tích hợp speaker verification gate} để cá nhân
          hóa hệ thống KWS.
\end{enumerate}

% ============================================================
\chapter{Cơ Sở Lý Thuyết}

\section{Tổng quan Keyword Spotting}

KWS biến bài toán phát hiện từ khóa trong luồng âm thanh liên tục thành bài
toán phân lớp thời gian thực thông qua \textbf{phương pháp cửa sổ}: luồng âm
thanh được chia thành các đoạn ngắn, chồng lấp (thường 1 giây), và mỗi đoạn
được phân lớp độc lập.

Pipeline KWS tiêu chuẩn gồm:
\begin{enumerate}
    \item \textbf{Tiền xử lý âm thanh:} Tín hiệu thô $\rightarrow$ trích
          xuất đặc trưng (MFCC)
    \item \textbf{Mã hóa đặc trưng:} Mạng neural ánh xạ đặc trưng vào không
          gian embedding
    \item \textbf{Phân lớp:} Xác định từ khóa nào (nếu có) xuất hiện
\end{enumerate}

\section{Trích xuất Đặc trưng Âm thanh: MFCC}
\label{sec:mfcc_vi}

MFCC (Mel-Frequency Cepstral Coefficients) là phương pháp biểu diễn đặc trưng
âm thanh phổ biến nhất cho KWS. Quy trình trích xuất gồm các bước:

\subsection{Tiền nhấn mạnh (Pre-emphasis)}

Khuếch đại thành phần tần số cao của tín hiệu:
\begin{equation}
    x[n] = y[n] - \alpha \cdot y[n-1], \quad \alpha = 0.97
\end{equation}

\subsection{Phân khung và Cửa sổ (Framing \& Windowing)}

Tín hiệu được chia thành các khung chồng lấp:
\begin{equation}
    x_m[n] = x[n + m \cdot H], \quad 0 \leq n < N, \quad 0 \leq m < M
\end{equation}

\textbf{Trong hệ thống của chúng tôi:}
\begin{itemize}
    \item Chiều dài khung $N = 640$ mẫu (40ms tại 16kHz)
    \item Bước nhảy $H = 320$ mẫu (20ms)
    \item Số khung $M \approx 49$ cho âm thanh 1 giây
\end{itemize}

Mỗi khung được làm mịn bằng cửa sổ Hamming:
\begin{equation}
    w[n] = 0.54 - 0.46 \cos\left(\frac{2\pi n}{N-1}\right)
\end{equation}

\subsection{Biến đổi Fourier Rời rạc}

Chuyển đổi từ miền thời gian sang miền tần số:
\begin{equation}
    X[k] = \sum_{n=0}^{N-1} x_m[n] \cdot e^{-j2\pi kn/N}
\end{equation}

Phổ công suất: $P[k] = |X[k]|^2 / N$

\subsection{Bộ lọc Mel}

Phổ công suất được lọc qua $F = 40$ bộ lọc tam giác trên thang Mel:
\begin{equation}
    S_m = \sum_{k=f_{m-1}}^{f_{m+1}} P[k] \cdot H_m[k]
\end{equation}

Công thức chuyển đổi sang thang Mel:
\begin{equation}
    \mathcal{M}(f) = 2595 \cdot \log_{10}\left(1 + \frac{f}{700}\right)
\end{equation}

\subsection{DCT và MFCC}

Áp dụng Biến đổi Cosine Rời rạc lên log Mel-spectrogram:
\begin{equation}
    C[k] = 2\beta \sum_{f=0}^{F-1} S_f \cdot \cos\left(\frac{\pi k(2f+1)}{2F}\right)
\end{equation}

\textbf{Đầu ra:} 40 hệ số MFCC được tính, chỉ giữ 10 đầu tiên (theo Rusci et al.).
Quy trình: MFCC(40) $(1, 40, 49)$ $\rightarrow$ narrow(10) $(1, 10, 49)$ $\rightarrow$
transpose $(1, 49, 10)$ = (channel, $T$, features). Tensor đầu vào DSCNN: $(B, 1, 49, 10)$.


\section{Mạng Tích chập Phân tách theo Chiều sâu (DSCNN)}
\label{sec:dscnn_vi}

DSCNN phân tách phép tích chập 3D tiêu chuẩn thành hai phép toán nhẹ hơn:

\subsection{Tích chập theo Chiều sâu (Depthwise Convolution)}

Áp dụng một bộ lọc duy nhất cho mỗi kênh đầu vào:
\begin{equation}
    Y^{(c)} = X^{(c)} * K^{(c)}
\end{equation}

\subsection{Tích chập theo Điểm (Pointwise Convolution)}

Kết hợp thông tin giữa các kênh bằng tích chập $1 \times 1$:
\begin{equation}
    Z_k = \sum_{c=1}^{C} Y^{(c)} * W_k^{(c)}
\end{equation}

\subsection{Kiến trúc DSCNN-L}

Mô hình DSCNN-L (276 kênh, theo Rusci et al.\ \texttt{model\_size\_info\_DSCNNL})
gồm 1 lớp conv ban đầu + 5 khối DS:

\begin{itemize}
    \item \textbf{Conv ban đầu:} Conv2d$(1, 276, \text{kernel}=(10,4), \text{stride}=(2,1))$ + BN + ReLU
    \item \textbf{Khối DS 1:} stride $(2,2)$; \textbf{Khối DS 2--5:} stride $(1,1)$
    \item \textbf{Khối DS 1--4:} DW + BN + ReLU + PW + BN + ReLU
    \item \textbf{Khối DS 5 (cuối):} DW vẫn dùng BN+ReLU; chỉ sau PW dùng LayerNorm(\texttt{elementwise\_affine=False})
    \item \textbf{Average Pooling} $\rightarrow$ Flatten $\rightarrow$ đầu ra $(B, 276)$
    \item \textbf{Không có Linear projection:} embedding\_dim = số kênh = 276
    \item \textbf{Chuẩn hóa L2:} áp dụng \emph{bên ngoài} model (\texttt{F.normalize} trong ReprModel)
\end{itemize}

\begin{equation}
    \hat{e} = \frac{e}{\|e\|_2}
\end{equation}

\textbf{Lưu ý quan trọng:} Chuẩn hóa L2 đảm bảo tất cả embedding nằm trên
siêu cầu đơn vị, giúp khoảng cách L2 tương đương khoảng cách cosine (sai khác
hằng số).


\section{Mạng Prototypical (Prototypical Networks)}
\label{sec:protonet_vi}

Mạng Prototypical học một không gian embedding nơi phân lớp được thực hiện
bằng cách tính khoảng cách đến các \textbf{prototype} của lớp.

\subsection{Ý tưởng Cốt lõi}

Cho mạng neural $f_\theta : \mathbb{R}^D \rightarrow \mathbb{R}^M$ (DSCNN
encoder). Với mỗi lớp $k$ có tập hỗ trợ $S_k$, \textbf{prototype} $c_k$
là trung bình embedding của các mẫu hỗ trợ:

\begin{equation}
    c_k = \frac{1}{|S_k|} \sum_{(x_i, y_i) \in S_k} f_\theta(x_i)
\end{equation}

\textbf{Cách hiểu trực quan:} Prototype là ``dấu vân tay đại diện'' của mỗi
lớp trong không gian embedding. Mẫu mới được phân lớp dựa trên prototype
gần nhất.

\subsection{Huấn luyện theo Tập (Episodic Training)}

Huấn luyện mô phỏng kịch bản few-shot qua các \textbf{episode}:

\begin{enumerate}
    \item Chọn ngẫu nhiên $N$ lớp từ tập huấn luyện
    \item Mỗi lớp: lấy $K$ mẫu hỗ trợ + $Q$ mẫu truy vấn
    \item Tính prototype từ tập hỗ trợ
    \item Phân lớp mẫu truy vấn và tính hàm mất mát
\end{enumerate}


\section{Hàm Mất mát Triplet (Triplet Loss)}
\label{sec:triplet_vi}

Hệ thống sử dụng \textbf{Triplet Loss} thay vì loss prototypical tiêu chuẩn:

\begin{equation}
    \mathcal{L}_{TL} = \frac{1}{N_t} \sum_{i=1}^{N_t}
    \max\left(0, \; d(x_i, x_i^+) - d(x_i, x_i^-) + m\right)
\end{equation}

trong đó:
\begin{itemize}
    \item $x_i$: mẫu \textbf{neo} (anchor)
    \item $x_i^+$: mẫu \textbf{dương} (positive -- cùng lớp với neo)
    \item $x_i^-$: mẫu \textbf{âm} (negative -- khác lớp)
    \item $m = 0.5$: \textbf{biên} (margin)
    \item $d(\cdot, \cdot)$: khoảng cách L2
\end{itemize}

\textbf{Trực giác:} Loss đẩy neo lại gần mẫu dương và ra xa mẫu âm ít nhất
một khoảng biên $m$. Khi ràng buộc đã thỏa mãn, loss bằng 0 (hành vi
hinge loss).

\textbf{Cấu hình huấn luyện:}
\begin{itemize}
    \item Batch: 20 mẫu $\times$ 80 lớp, margin $m = 0.5$
    \item Optimizer: Adam, lr = 0.001, StepLR(step\_size=20, $\gamma_{\text{decay}}=0.5$)
    \item 40 epoch $\times$ 400 episode
    \item Dữ liệu: MSWC English (top 500 từ)
    \item Augmentation: Nhiễu DEMAND, $p = 0.95$, SNR cố định $= 5$\,dB
\end{itemize}


\section{Phân lớp Open-Set}
\label{sec:openset_vi}

Phân lớp open-set mở rộng phân lớp tiêu chuẩn bằng khả năng \textbf{từ chối
mẫu không xác định} --- mẫu không thuộc bất kỳ lớp từ khóa nào đã đăng ký.

\subsection{Phương pháp Khoảng cách L2 Trực tiếp (Đề xuất của chúng tôi)}

Thay vì chuẩn hóa xác suất, sử dụng khoảng cách L2 trực tiếp:

\begin{equation}
    d_k = \|f_\theta(x_q) - c_k\|_2, \quad k = 1, \ldots, N
\end{equation}

\begin{equation}
    \hat{y} = \begin{cases}
        \displaystyle\arg\min_k d_k & \text{nếu } \min_k d_k \leq \gamma \\
        \text{unknown} & \text{ngược lại}
    \end{cases}
\end{equation}

trong đó $\gamma$ là ngưỡng chấp nhận (bán kính).

\textbf{Ưu điểm so với phương pháp xác suất:}
\begin{itemize}
    \item Ngưỡng $\gamma$ có ý nghĩa hình học (bán kính chấp nhận quanh prototype)
    \item Không phụ thuộc vào prototype ``unknown'' ảo
    \item Hiệu năng tốt hơn: AUC 0.95 vs 0.94, ACC@5\%FAR 0.81 vs 0.76
          (kết quả tham chiếu)
\end{itemize}


\section{Phát hiện Hoạt động Giọng nói (VAD)}

VAD xác định đoạn âm thanh nào chứa tiếng người nói. Hệ thống sử dụng
\textbf{Silero VAD} --- mô hình pre-trained nhẹ, xử lý chunk 32ms (512 mẫu
tại 16kHz), đầu ra xác suất tiếng nói $p \in [0, 1]$.

Tiếng nói được phát hiện khi $p > \tau_{vad}$ (mặc định $\tau_{vad} = 0.5$).


\section{Xác thực Người nói với ECAPA-TDNN}

ECAPA-TDNN tạo biểu diễn người nói 192 chiều. Xác thực bằng cosine similarity:

\begin{equation}
    \text{similarity}(e_{owner}, e_{query}) =
    \frac{e_{owner} \cdot e_{query}}
         {\|e_{owner}\|_2 \cdot \|e_{query}\|_2}
\end{equation}


% ============================================================
\chapter{Thiết Kế Hệ Thống}

\section{Tổng quan Kiến trúc}

Hệ thống theo kiến trúc pipeline module hóa với ba tầng phân rõ:

\begin{itemize}
    \item \textcolor{corecolor}{\textbf{CORE:}} KWS engine (MFCC + DSCNN +
          OpenNCM) --- \textit{bắt buộc}
    \item \textcolor{ext1color}{\textbf{EXT-1:}} Kháng nhiễu (denoising) ---
          \textit{mở rộng 1}
    \item \textcolor{ext2color}{\textbf{EXT-2:}} Streaming inference (VAD +
          sliding window) --- \textit{mở rộng 2}
    \item \textcolor{optcolor}{\textbf{OPTIONAL:}} Speaker verification ---
          \textit{tùy chọn}
\end{itemize}

\textbf{Nguyên tắc thiết kế:} Mỗi module tùy chọn có thể bằng \texttt{None}
trong pipeline. Nếu module lỗi, hệ thống tiếp tục hoạt động bình thường
(fail-open).


\section{Quy trình Ba Giai đoạn}

\subsection{Giai đoạn A: Huấn luyện (Offline, Một lần)}

DSCNN encoder $f_\theta$ được huấn luyện trên MSWC English bằng Triplet Loss.
Sau huấn luyện, trọng số encoder được \textbf{đóng băng} (frozen).

\subsection{Giai đoạn B: Đăng ký (Few-shot, Mỗi người dùng)}

Người dùng cung cấp $K$ mẫu âm thanh (thường $K = 3$--$5$) cho mỗi từ khóa.
Mẫu được đưa qua encoder đóng băng để tính prototype:

\begin{equation}
    c_k = \frac{1}{K} \sum_{i=1}^{K} f_\theta(x_i^{(k)})
\end{equation}

\subsection{Giai đoạn C: Suy luận}

Mỗi mẫu truy vấn được tính khoảng cách L2 đến tất cả prototype và áp dụng
quy tắc quyết định open-set.


\section{Mở rộng 1: Kháng Nhiễu}

\subsection{Khi huấn luyện: Noise Augmentation}

Trộn âm thanh sạch với nhiễu nền từ DEMAND (SNR cố định 5\,dB, theo \texttt{noise\_snr=5}):
\begin{equation}
    x_{noisy} = x_{clean} + \alpha \cdot x_{noise}, \quad
    \alpha = \frac{\text{RMS}(x_{clean})}{\text{RMS}(x_{noise}) \cdot 10^{\text{SNR}/20}}
\end{equation}

\subsection{Khi suy luận: Khử nhiễu (Denoising)}

Hai phương pháp:
\begin{itemize}
    \item \textbf{Spectral Gating} (nhanh, CPU): Khử nhiễu thống kê
    \item \textbf{SepformerSeparation} (chính xác, GPU): Khử nhiễu neural
\end{itemize}

\textbf{Kế hoạch benchmark:} So sánh accuracy KWS tại SNR 0, 5, 10, 20 dB
và sạch, trước và sau khử nhiễu.


\section{Mở rộng 2: Suy luận Streaming}

\subsection{Phát hiện Cửa sổ Trượt}

\begin{itemize}
    \item Kích thước cửa sổ: 16,000 mẫu (1 giây tại 16kHz)
    \item Bước trượt: 8,000 mẫu (0.5 giây, chồng lấp 50\%)
    \item Buffer: Ring buffer (\texttt{collections.deque})
\end{itemize}

\subsection{Tích hợp VAD}

Silero VAD lọc trước các chunk âm thanh (512 mẫu = 32ms). KWS chỉ được kích
hoạt khi phát hiện tiếng nói, tiết kiệm tài nguyên tính toán.

\subsection{Ngân sách Độ trễ}

\begin{table}[h]
\centering
\caption{Độ trễ mục tiêu cho mỗi thành phần pipeline}
\begin{tabular}{lr}
\toprule
\textbf{Thành phần} & \textbf{Mục tiêu} \\
\midrule
VAD mỗi chunk (32ms) & $< 5$ ms \\
Khử nhiễu (1 giây) & $< 30$ ms \\
Trích xuất MFCC & $< 5$ ms \\
Suy luận DSCNN & $< 10$ ms \\
Phân lớp & $< 1$ ms \\
Xác thực người nói & $< 20$ ms \\
\midrule
\textbf{Tổng cộng} & $\mathbf{< 71}$ \textbf{ms} \\
\bottomrule
\end{tabular}
\end{table}


\section{Tùy chọn: Speaker Verification Gate}

Sau khi KWS phát hiện từ khóa, gate xác thực người nói kiểm tra danh tính:

\begin{equation}
    \hat{y}_{final} = \begin{cases}
        \hat{y}_{kws} & \text{nếu } \text{sim}(e_{owner}, e_{query}) > \tau_{spk} \\
        \text{từ chối} & \text{ngược lại}
    \end{cases}
\end{equation}

Module này hoàn toàn \textbf{tùy chọn} --- hệ thống hoạt động bình thường
khi không bật.


% ============================================================
\chapter{Kế Hoạch Đánh Giá}

\section{Bộ Dữ liệu}

\subsection{Google Speech Commands v2 (GSC)}

\begin{itemize}
    \item 105,829 phát ngôn của 35 từ tiếng Anh bởi 2,618 người nói
    \item File WAV 1 giây, tần số lấy mẫu 16kHz
    \item Dùng cho \textbf{đánh giá} (enrollment + evaluation)
\end{itemize}

\textbf{Phân vùng từ (Fixed protocol):}
\begin{itemize}
    \item \textbf{GSC+ (10 positive):} yes, no, up, down, left, right, on,
          off, stop, go
    \item \textbf{GSC$-$ (20 negative):} các từ còn lại (lớp unknown)
    \item \textbf{Loại trừ (5):} backward, forward, visual, follow, learn
\end{itemize}

\subsection{MSWC (Multilingual Spoken Words Corpus)}

\begin{itemize}
    \item $\sim$22.5 triệu phát ngôn, 50 ngôn ngữ
    \item Dùng cho \textbf{huấn luyện} DSCNN encoder
    \item Subset English: top 500 từ theo số phát ngôn
    \item Subset đánh giá: 263 từ (min 1000 phát ngôn, loại trừ 500 từ huấn luyện)
\end{itemize}

\section{Chỉ số Đánh giá}

\begin{itemize}
    \item \textbf{Đường cong DET trung bình:} FAR vs FRR trung bình trên
          tất cả lớp từ khóa
    \item \textbf{AUC:} Diện tích dưới đường cong ROC (keyword vs non-keyword)
    \item \textbf{ACC@5\%FAR:} Accuracy tại điểm FAR = 5\%
\end{itemize}

\section{Giao thức Đánh giá}

\begin{table}[h]
\centering
\caption{Giao thức đánh giá theo mức ưu tiên}
\begin{tabular}{llll}
\toprule
\textbf{Ưu tiên} & \textbf{Giao thức} & \textbf{Positive} & \textbf{Negative} \\
\midrule
Bắt buộc & GSC Fixed & 10 từ cố định & 20 từ cố định \\
Bắt buộc & GSC Randomized & 10 ngẫu nhiên/lần & 20 ngẫu nhiên/lần \\
Mở rộng & MSWC Randomized & 5 ngẫu nhiên & 50 ngẫu nhiên \\
\bottomrule
\end{tabular}
\end{table}

\section{Đảm bảo Không Data Leakage}

\begin{itemize}
    \item Encoder huấn luyện trên MSWC English (top 500 từ)
    \item Đánh giá trên GSC (35 từ hoàn toàn khác)
    \item Enrollment, calibration, và test sử dụng split tách biệt
    \item Query test không bao giờ dùng để tạo prototype hoặc chọn threshold
\end{itemize}


% ============================================================
\chapter{Kế Hoạch Triển Khai}

\section{Lộ trình}

\begin{table}[h]
\centering
\caption{Lộ trình triển khai với milestone và mức ưu tiên}
\begin{tabular}{clll}
\toprule
\textbf{Tuần} & \textbf{Giai đoạn} & \textbf{Sản phẩm} & \textbf{Ưu tiên} \\
\midrule
1 & Setup + Pipeline & MFCC $\to$ DSCNN $\to$ embedding &
    \textcolor{corecolor}{\textbf{CORE}} \\
2 & Huấn luyện + Classifier & Model + enroll/predict &
    \textcolor{corecolor}{\textbf{CORE}} \\
3 & Đánh giá & Bảng kết quả + DET curves &
    \textcolor{corecolor}{\textbf{CORE}} \\
\midrule
\multicolumn{4}{c}{\textit{--- Điểm dừng an toàn: đủ để bảo vệ đồ án ---}} \\
\midrule
4 & Khử nhiễu & Bảng benchmark SNR &
    \textcolor{ext1color}{\textbf{EXT-1}} \\
5 & Streaming + VAD & Sliding window demo &
    \textcolor{ext2color}{\textbf{EXT-2}} \\
6 & Speaker Verify & Demo ``chỉ chủ nhân'' &
    \textcolor{optcolor}{\textbf{TÙY CHỌN}} \\
7 & Demo & Gradio web app 3 tab &
    DEMO \\
8 & Tài liệu & Báo cáo + slides &
    DOCS \\
\bottomrule
\end{tabular}
\end{table}

\section{Chiến lược Cắt Lỗ}

Nếu chậm tiến độ, cắt theo thứ tự từ dưới lên:
\begin{enumerate}
    \item \textbf{Cắt đầu tiên:} Speaker verification
    \item \textbf{Cắt tiếp:} Streaming realtime mic (giữ sliding window offline)
    \item \textbf{Cắt cuối:} Denoising (giữ noise augmentation khi train)
    \item \textbf{KHÔNG BAO GIỜ cắt:} Core KWS + Evaluation + Offline demo
\end{enumerate}


% ============================================================
\chapter{Kết Quả Kỳ Vọng}

Dựa trên kết quả tham chiếu từ luận văn gốc, chúng tôi kỳ vọng kết quả
cùng xu hướng tương tự:

\begin{table}[h]
\centering
\caption{Mốc hiệu năng tham chiếu}
\begin{tabular}{lccc}
\toprule
\textbf{Phương pháp} & \textbf{AUC} $\uparrow$ &
\textbf{ACC@5\%FAR} $\uparrow$ & \textbf{FRR@5\%FAR} $\downarrow$ \\
\midrule
Dựa trên xác suất & 0.94 & 0.76 & 0.21 \\
Khoảng cách L2 trực tiếp & 0.95 & 0.81 & 0.17 \\
\bottomrule
\end{tabular}
\end{table}

\textbf{Dự kiến phát hiện:}
\begin{enumerate}
    \item Phương pháp L2 trực tiếp vượt trội so với xác suất tại mọi
          điểm hoạt động
    \item Giao thức chọn từ ngẫu nhiên cho kết quả thực tế hơn (thấp hơn)
          so với cố định
    \item Khử nhiễu cải thiện accuracy rõ rệt nhất tại SNR thấp (0--5 dB)
    \item Hiệu năng few-shot ổn định từ 5 mẫu enrollment trở lên
    \item Streaming có thể giảm nhẹ accuracy so với offline do hiệu ứng
          biên cửa sổ
\end{enumerate}


% ============================================================
\begin{thebibliography}{9}

\bibitem{lopez2021deep}
I.~L{\'o}pez-Espejo et al.,
``Deep spoken keyword spotting: An overview,''
\textit{IEEE Access}, vol.~10, pp.~4169--4199, 2021.

\bibitem{snell2017prototypical}
J.~Snell, K.~Swersky, and R.~Zemel,
``Prototypical networks for few-shot learning,''
\textit{NeurIPS}, vol.~30, 2017.

\bibitem{rusci2023}
M.~Rusci and T.~Tuytelaars,
``Few-shot open-set learning for on-device customization of keyword spotting
systems,''
\textit{arXiv:2306.02161}, 2023.

\bibitem{zhang2017hello}
Y.~Zhang et al.,
``Hello edge: Keyword spotting on microcontrollers,''
\textit{arXiv:1711.07128}, 2017.

\bibitem{binh2025}
P.~T.~Binh,
``Enhanced evaluation protocols for few-shot open-set keyword spotting systems,''
Luận văn Thạc sĩ, USTH, 2025.

\bibitem{warden2018}
P.~Warden,
``Speech commands: A dataset for limited-vocabulary speech recognition,''
\textit{arXiv:1804.03209}, 2018.

\bibitem{mazumder2021mswc}
M.~Mazumder et al.,
``Multilingual spoken words corpus,''
\textit{NeurIPS Datasets and Benchmarks}, 2021.

\bibitem{speechbrain}
SpeechBrain Team,
``SpeechBrain: A general-purpose speech toolkit,'' 2021.

\end{thebibliography}

\end{document}
